\documentclass[11pt,a4paper]{article}

% ---------------------------------------------------------------- packages
\usepackage[T1]{fontenc}
\usepackage{amsmath,amssymb,amsthm}
\usepackage{newpxtext,newpxmath}
\usepackage[scaled=0.92]{inconsolata}
\usepackage[a4paper,margin=2.6cm]{geometry}
\usepackage{microtype}
\usepackage{mathtools}
\usepackage{graphicx}
\usepackage[font=small,labelfont=bf,labelsep=period]{caption}
\usepackage{subcaption}
\usepackage{booktabs}
\usepackage{tabularx}
\usepackage{longtable}
\usepackage{enumitem}
\usepackage{xcolor}
\usepackage{thmtools}
\usepackage[colorlinks=true,linkcolor=blue!50!black,citecolor=blue!50!black,urlcolor=blue!50!black]{hyperref}
\usepackage[capitalize,nameinlink]{cleveref}
\usepackage{fancyvrb}
\usepackage{newunicodechar}
% Unicode in quoted Lean code
\newunicodechar{ℝ}{\ensuremath{\mathbb{R}}}
\newunicodechar{↔}{\ensuremath{\leftrightarrow}}
\newunicodechar{√}{\ensuremath{\surd}}
\newunicodechar{φ}{\ensuremath{\varphi}}
\newunicodechar{₁}{\textsubscript{1}}
\newunicodechar{₂}{\textsubscript{2}}
\newunicodechar{₃}{\textsubscript{3}}

% ---------------------------------------------------------------- theorems
\declaretheorem[name=Theorem]{theorem}
\declaretheorem[name=Proposition,sibling=theorem]{proposition}
\declaretheorem[name=Lemma,sibling=theorem]{lemma}
\declaretheorem[name=Remark,style=remark,numbered=no]{remark}

% ---------------------------------------------------------------- notation
\newcommand{\C}{\mathbb{C}}
\newcommand{\R}{\mathbb{R}}
\newcommand{\Z}{\mathbb{Z}}
\newcommand{\Q}{\mathbb{Q}}
\newcommand{\GG}{\mathrm{GG}}
\newcommand{\rc}{r_{\mathrm c}}
\newcommand{\rs}{r_{\ast}}
\newcommand{\Lam}{\Lambda}
\newcommand{\Or}{\mathcal{O}}
\newcommand{\cls}{\operatorname{cl}}
\newcommand{\Win}{W}
\newcommand{\bm}{\bar m}
\newcommand{\sym}{\operatorname{Sym}}
\newcommand{\re}{\operatorname{Re}}
\newcommand{\im}{\operatorname{Im}}
% Lean declarations
\newcommand{\leanname}[1]{\texttt{\small #1}}
% Items the authors still have to supply.
\newcommand{\placeholder}[1]{\textcolor{red!70!black}{[#1]}}

\title{\textbf{The critical radius of the five-fold two-disk puzzle}}
\author{\placeholder{Authors}}
\date{\placeholder{Date}}

\begin{document}
\maketitle

\begin{abstract}
\noindent
Two closed disks of radius $r$, centred at $-1$ and $+1$ and each turned clockwise by $72^\circ$,
generate a group $\GG_5(r)$ of piecewise rotations of the plane: the group of an idealised two-disk
circle puzzle. Hearn, Kretschmer, Rokicki, Streeter and Vergo proved that $\GG_5(\sqrt{3+\varphi})$ is
infinite, where $\varphi$ is the golden ratio, and their numerical estimate of the radius at which
the group becomes infinite agrees with $\sqrt{3+\varphi}$ to six decimals. We prove that
$\GG_5(r)$ is finite for every $r<\sqrt{3+\varphi}$. Hence the group is finite exactly when
$r<\sqrt{3+\varphi}$, and the critical radius is $\sqrt{3+\varphi}=2.1489611417\ldots$ The finiteness
proof reads every position of a point off a lattice point of $\Z[e^{2\pi i/5}]$ and follows an
integer-valued \emph{level} that can only move while the point lies in the lens where the two disks
overlap; below $\sqrt{3+\varphi}$ the lens is too thin for the level to pass certain values. Both
halves are formalised in Lean~4 with Mathlib, and the formal statement is checked against an
independent, elementary restatement.
\end{abstract}

% =================================================================================================
\section{The puzzle and the number}\label{sec:intro}

A two-disk circle puzzle such as Gizmo Gears consists of two overlapping disks that turn
independently; a turn carries along everything in the turning disk, including the part of the other
disk that it covers. Whether the puzzle can be cut into finitely many pieces that the turns merely
permute depends on the overlap: with a small overlap it can, but past a critical overlap the turns
cut ever smaller pieces and the puzzle \emph{jumbles}. Hearn, Kretschmer, Rokicki, Streeter and
Vergo~\cite{hkrsv} turned this into a question about groups. Let $r\in\R$ and
\[
  D_- = \{z\in\C : |z+1|\le r\},\qquad D_+ = \{z\in\C : |z-1|\le r\},\qquad \zeta = e^{2\pi i/5}.
\]
The \emph{turns} are the maps
\[
  a(z)=\begin{cases}\zeta^{-1}(z+1)-1, & z\in D_-,\\ z, & \text{otherwise,}\end{cases}\qquad
  b(z)=\begin{cases}\zeta^{-1}(z-1)+1, & z\in D_+,\\ z, & \text{otherwise.}\end{cases}
\]
Each rotates its closed disk clockwise by $72^\circ$ and fixes every other point, so it is a
bijection of~$\C$. As on a puzzle, words are read from left to right: $ab$ means ``first $a$, then
$b$''. The group
\[
  \GG_5(r) = \langle a, b\rangle \le \sym(\C)
\]
is the \emph{two-disk group}; its elements are piecewise rotations, not rotations. The
\emph{orbit} $\Or_r(p)=\{g(p): g\in\GG_5(r)\}$ is the set of places a sticker at $p$ can reach.
Enlarging the disks can only enlarge orbits, and a finitely generated group of permutations with
uniformly bounded orbits is finite (\cref{prop:mono,prop:bounded} below); so the radii at which
$\GG_5(r)$ is finite form an initial segment of $\R$, and its supremum
\[
  \rc(5) = \sup\{r\in\R : \GG_5(r)\text{ is finite}\}
\]
is the \emph{critical radius}. For $0<r<1$ the disks are disjoint and $\GG_5(r)\cong C_5\times C_5$~\cite{hkrsv};
as $r$ grows the orbits grow (\cref{fig:orbit}), and at the critical radius they become infinite.

\begin{figure}[t]
  \centering
  \includegraphics[width=0.62\textwidth]{figures/b4_monotone.png}
  \caption{The orbit of the point $0$ under $\GG_5(2.13)$: $2\,543$ points, drawn with the two
    disks. Slightly below the critical radius the orbits are finite but large; at
    $\rs=\sqrt{3+\varphi}$ this orbit becomes infinite.}
  \label{fig:orbit}
\end{figure}

Hearn et al.\ proved that $\GG_5$ is infinite at $r=\sqrt{3+\varphi}$~\cite[Thm.~2]{hkrsv}, where
$\varphi=(1+\sqrt5)/2$, so that $\rc(5)\le\sqrt{3+\varphi}$, and their numerical estimate
$\rc(5)\approx 2.148961$ agreed with $\sqrt{3+\varphi}=2.1489611\ldots$~\cite[Table~1]{hkrsv}. That
the group is finite for every smaller radius was not known. We prove it.

% =================================================================================================
\section{The result}\label{sec:result}

Throughout, $\rs=\sqrt{3+\varphi}=2.1489611417\ldots$; it is the largest root of $x^4-7x^2+11$
(its conjugate $\sqrt{4-\varphi}$ is the value found in~\cite{hkrsv} for $\GG_{10}$).

\begin{theorem}\label{thm:main}
$\GG_5(r)$ is finite if and only if $r<\rs$. In particular $\rc(5)=\rs=\sqrt{3+\varphi}$. More precisely:
\begin{enumerate}[label=(\alph*),nosep]
  \item at $r=\rs$ the orbit of $0$ is infinite, so $\GG_5(r)$ is infinite for every $r\ge\rs$;
  \item for every $r<\rs$ there is $N(r)$ such that every orbit of $\GG_5(r)$ has at most $N(r)$
        points, so $\GG_5(r)$ is finite.
\end{enumerate}
\end{theorem}

Both statements of the first sentence are formal theorems (\cref{sec:formal}):
\begin{Verbatim}[fontsize=\small,xleftmargin=1.5em]
theorem finite_GG_five_iff (r : ℝ) : Finite (GG 5 r) ↔ r < √(3 + φ)
theorem criticalRadius_five : criticalRadius 5 = √(3 + φ)
\end{Verbatim}
where \texttt{criticalRadius n} is \texttt{sSup \{r | Finite (GG n r)\}}. By (a) and (b) this set is
exactly $(-\infty,\rs)$, so the supremum is a genuine one; in particular the group is infinite at
the critical radius itself. Two general facts connect orbits, groups and radii.

\begin{proposition}[monotonicity]\label{prop:mono}
If $r\le r'$ then $\Or_r(p)\subseteq\Or_{r'}(p)$ for every $p\in\C$.
\end{proposition}

\begin{proof}
Apply a word to $p$ letter by letter. A letter that moves the current point at radius $r$ does so
because the point lies in its radius-$r$ disk, hence in the concentric radius-$r'$ disk, where the
radius-$r'$ letter moves it in the same way. A letter that fixes the current point is dropped.
\end{proof}

Note that $\GG_5(r)$ is not a subgroup of $\GG_5(r')$: the comparison is made point by point.

\begin{proposition}[bounded orbits]\label{prop:bounded}
Let $G$ be a group of permutations of a set $X$ generated by a finite set. If every $G$-orbit has
at most $N$ points, then $G$ is finite.
\end{proposition}

\begin{proof}
For $x\in X$, number the orbit of $x$ and let $\rho_x\colon G\to S_N$ be the action on it. Each
$\rho_x$ is determined by the images of the finitely many generators, so only finitely many
homomorphisms $\rho_x$ occur. An element of $G$ in the kernel of every $\rho_x$ fixes every point,
hence is the identity; so $g\mapsto(\rho_x(g))_x$ embeds $G$ into a finite product of copies of~$S_N$.
\end{proof}

Hence, if $\GG_5(r')$ is finite and $r\le r'$, then $\GG_5(r)$ is finite (its orbits lie in those of
$\GG_5(r')$, which have at most $|\GG_5(r')|$ points); and an infinite orbit at radius $r$ makes
$\GG_5(r')$ infinite for every $r'\ge r$. \Cref{thm:main} therefore follows from its two parts, which
are proved in \cref{sec:upper,sec:lower}.

% =================================================================================================
\section{Why it jumbles at \texorpdfstring{$\sqrt{3+\varphi}$}{sqrt(3+phi)}}\label{sec:upper}

This is the construction of~\cite[Thm.~2]{hkrsv}; we recall it briefly. Let $r=\rs$ and
\[
  E=\zeta-\zeta^2,\qquad F=E/\varphi=1-\zeta+\zeta^2-\zeta^3,\qquad G=2F-E,
\]
and write $X'=-X$. Then $|E+1|^2=3+\varphi$: the point $E$ lies on the rim of $D_-$, and $E'$
(by the symmetry $z\mapsto-z$) on the rim of $D_+$. The chord $E'E$ passes through $0$ and spans the
lens $D_-\cap D_+$ from rim to rim (\cref{fig:setting}). The points
$E',F',G',G,F,E$ lie on it in this order and cut it into three pieces
\[
  I_1=[E',F'],\qquad I_2=[F',G'],\qquad I_3=[G',E].
\]

\begin{figure}[t]
  \centering
  \includegraphics[width=0.56\textwidth]{figures/a1_setting.png}
  \caption{The two disks at $r=\rs$, the lens, and the chord $E'E$ with its three pieces $I_1$
    (green), $I_2$ (magenta) and $I_3$ (yellow).}
  \label{fig:setting}
\end{figure}

\begin{lemma}[three words]\label{lem:words}
At $r=\rs$, with $A=a^{-1}$ and $B=b^{-1}$,
\begin{enumerate}[label=(\roman*),nosep]
  \item $AABAB=a^{-2}b^{-1}a^{-1}b^{-1}$ maps $I_1$ onto $[G,F]$, translating it by $2F$;
  \item $ababb=abab^2$ maps $I_2$ onto $[F,E]$, translating it by $2F$;
  \item $abaBAB=abab^{-1}a^{-1}b^{-1}$ maps $I_3$ onto $[E',G]$, translating it by $2F-2E$.
\end{enumerate}
\end{lemma}

\begin{proof}
A word acts on a segment as the composition of the rotations of its letters as long as every
intermediate image of the segment lies in the disk of the next letter. Disks are convex, so it is
enough to check the two endpoints of each intermediate segment. These endpoints are explicit
elements of $\Z[\zeta]$, and for $x\in\Z[\zeta]$ the number $|x\pm1|^2$ lies in $\Z[\varphi]$, so
each check is an exact comparison $m+n\varphi\le 3+\varphi$ with $m,n\in\Z$. In each word the
rotation angles add up to $0$, so the composition is a translation, and the endpoints go where
stated (\cref{fig:words}).
\end{proof}

\begin{figure}[t]
  \centering
  \begin{subfigure}{0.325\textwidth}\includegraphics[width=\linewidth]{figures/a2_word1.png}
    \caption{$AABAB$ on $I_1$}\end{subfigure}\hfill
  \begin{subfigure}{0.325\textwidth}\includegraphics[width=\linewidth]{figures/a3_word2.png}
    \caption{$ababb$ on $I_2$}\end{subfigure}\hfill
  \begin{subfigure}{0.325\textwidth}\includegraphics[width=\linewidth]{figures/a4_word3.png}
    \caption{$abaBAB$ on $I_3$}\end{subfigure}
  \caption{The three words, turn by turn: faint copies show the intermediate positions of the
    piece and arcs the paths of its midpoint; the arrow is the resulting translation. Red dots mark
    endpoints lying exactly on the rim of the disk about to turn.}
  \label{fig:words}
\end{figure}

\begin{remark}
Several of these comparisons are equalities: the endpoints touch the rims, at $\pm E$, at
$\pm\bar E$ and at the two tips $\pm(\zeta-\zeta^{-1})=\pm i\sqrt{2+\varphi}$ of the lens (the red
dots in \cref{fig:words}). For any smaller radius the same words would split the pieces.
\end{remark}

\begin{proposition}\label{prop:upper}
At $r=\rs$ the orbit of $0$ is infinite. Hence $\GG_5(r)$ is infinite for every $r\ge\rs$.
\end{proposition}

\begin{proof}
Parametrise the chord by $z=sE$, $s\in[-1,1)$. By \cref{lem:words} the three words realise the map
\[
  s\longmapsto\begin{cases} s+2/\varphi, & s<3-2\varphi,\\ s+2/\varphi-2, & s\ge 3-2\varphi,\end{cases}
\]
the rotation by $2/\varphi$ of the circle $[-1,1)$ of length $2$: the three pieces are exchanged,
$I_1\cup I_2$ and $I_3$ having lengths in the ratio $1:\varphi$ (\cref{fig:exchange}). Its rotation
number $1/\varphi$ is irrational, so the iterates of $s=0$ are pairwise distinct, and each of them is
the image of $0$ under an element of the group. By \cref{prop:mono} the orbit of $0$ is infinite at
every $r\ge\rs$, and an infinite orbit makes the group infinite.
\end{proof}

\begin{figure}[t]
  \centering
  \includegraphics[width=0.66\textwidth]{figures/a5_exchange.png}
  \caption{Top: the exchange of the three pieces on the chord, $I_1I_2I_3\mapsto I_3I_1I_2$. Bottom:
    gluing $E'$ to $E$ turns it into a rotation of a circle by $1/\varphi$ of a turn. The marks
    outside the circle (and above the chord) are the first $300$ points of the orbit of $0$,
    coloured by the piece they lie in; the orbit never closes.}
  \label{fig:exchange}
\end{figure}

% =================================================================================================
\section{Why it does not jumble below \texorpdfstring{$\sqrt{3+\varphi}$}{sqrt(3+phi)}}\label{sec:lower}

Fix $r<\rs$ and a point $p$. If $p$ lies in neither disk, its orbit is $\{p\}$; otherwise
$|p|\le 1+r$. Write $\Lam=\Z[\zeta]$ and $\mu_5=\{1,\zeta,\dots,\zeta^4\}$.

\subsection{The point's-eye view}

An element $g$ of the group acts on $p$ through the rotations of those letters that actually move
the current point, so $g(p)=T(p)$ for an isometry $T(z)=mz+\tau$ with $m\in\mu_5$ and
$\tau\in\Lam$. Each of the four rotations maps $-1+2\Lam$ onto itself; hence
$T^{-1}(-1)=2\nu-1$ with $\nu\in\Lam$, and
\begin{equation}\label{eq:normal}
  g(p) = m\,(p+1-2\nu)-1 .
\end{equation}
Read from the point's side, $T^{-1}$ carries the centres $-1$ and $+1$ of the disks to
$2\nu-1$ and $2(\nu+\bm)-1$: relative to the point the puzzle has moved, and its two centres now sit
at lattice points. We call $\nu$ the \emph{pivot} and $\nu+\bm$ the \emph{tip}; the segment joining
them is a \emph{dumbbell} of length~$1$. The next letter $a^{\pm1}$ moves $g(p)$ if and only if the
pivot lies in the \emph{window}
\[
  \Win_p=\{v\in\Lam : |p+1-2v|\le r\},
\]
and then it turns the dumbbell about the pivot ($m\mapsto\zeta^{\mp1}m$, $\nu$ fixed); the letter
$b^{\pm1}$ needs the tip in the window and turns the dumbbell about the tip
($m\mapsto\zeta^{\mp1}m$, $\nu\mapsto\nu+\bm-\zeta^{\pm1}\bm$). So an orbit is a walk of a dumbbell on the lattice that may
turn about an end only when that end lies in the window (\cref{fig:lattice}); both ends lie in the
window exactly when $p$ lies in the lens of the two disks about $2\nu-1$ and $2(\nu+\bm)-1$.

\begin{figure}[t]
  \centering
  \includegraphics[width=0.5\textwidth]{figures/b1_lattice.png}
  \caption{The point's-eye view of the orbit of a point $p$ at $r=2$ ($143$ positions). Each
    position is a pair of disk centres $2\nu-1$ (blue) and $2(\nu+\bm)-1$ (orange) joined by a
    segment; grey centres lie outside the window (dashed circle of radius $r$ about $p$) and cannot
    be turned about.}
  \label{fig:lattice}
\end{figure}

The lattice $\Lam$ has rank $4$ and is dense in $\C$, so the window contains infinitely many lattice
points and the walk could, a priori, wander forever inside it. What must stay bounded are the two
remaining coordinates, those of the Galois conjugate $\zeta\mapsto\zeta^2$. The rest of the proof
bounds them.

\paragraph{Types.}
The ideal $\mathfrak p=(1-\zeta)$ is a prime of $\Lam$ of norm $5$, and $\zeta\equiv1\pmod{\mathfrak p}$.
Write $\cls(v)\in\Z/5$ for the class of $v$ modulo $\mathfrak p$. Starting from the dumbbell
$(\nu,\nu+\bm)=(0,1)$ of the identity, induction along the word shows that the pivot always has
class $0$ and the tip class $1$. So the walk is \emph{typed}: from a point of class $0$ it adds a
fifth root of unity, from a point of class $1$ it subtracts one. Call $V\subseteq\Lam$
\emph{closed} if, for all $v\in V$ and $e\in\mu_5$,
\[
  \cls(v)=0,\ v+e\in\Win_p\ \Longrightarrow\ v+e\in V,\qquad
  \cls(v)=1,\ v-e\in\Win_p\ \Longrightarrow\ v-e\in V.
\]

\begin{lemma}[orbit invariant]\label{lem:invariant}
Let $V$ be closed, and suppose that each of $0,1$ lying in $\Win_p$ belongs to $V$ and that at
least one of them does. Then every point of $\Or_r(p)$ has the form \eqref{eq:normal} with
$m\in\mu_5$ and $\nu\in V$ or $\nu+\bm\in V$.
\end{lemma}

\begin{proof}
One shows, letter by letter, that every reachable state $(m,\nu)$ keeps the property of the initial
one: each end of the dumbbell that lies in $\Win_p$ is in $V$, and some end is in $V$. A turn happens
about an end in $\Win_p$, which is in $V$; the moved end, if it lands in $\Win_p$, is reached from it
by a typed step, so it is in $V$ as well.
\end{proof}

\subsection{A level}

Let $u=1-\zeta+\zeta^2$, so $|u|=2-\varphi$, and for $v\in\Lam$ let
\[
  \psi(v)=2\langle u,v\rangle=\bar uv+u\bar v\in\Z[\varphi],\qquad \psi(v)=A(v)+B(v)\,\varphi\quad(A,B\in\Z).
\]
The \emph{level} of $v$ is the integer $\ell(v)=-A(v)-3B(v)$, and we put $\delta=3-\varphi$. Both
$\psi$ and $\ell$ are additive.

\begin{lemma}[the level]\label{lem:level}
\begin{enumerate}[label=(\roman*),nosep]
  \item $\psi(v)+\ell(v)=-\delta\,B(v)$; so on each level $\ell$ the values of $\psi$ form the
        progression $-\ell+\delta\Z$.
  \item For $e\in\mu_5$, $\ell(e)\in\{3,-2\}$, and $\psi(e)=3-2\varphi$ whenever $\ell(e)=3$ (which
        happens for $e=1$ and $e=\zeta^2$).
  \item $\ell(v)\equiv 3\cls(v)\pmod 5$.
\end{enumerate}
\end{lemma}

\begin{proof}
(i) holds by definition; (ii) is a table of five values; (iii) is a congruence between integer
linear forms in the coordinates of $v$.
\end{proof}

So pivots sit on levels in $5\Z$ and tips on levels in $3+5\Z$, and a typed step changes the level
by $+3$ or $-2$ (from a pivot) or by $-3$ or $+2$ (from a tip). Consequently the walk can rise above
a level $5q$ only by an \emph{active} step: from a pivot $v$ on level $5q$ to $v+e$ with $\ell(e)=3$.
Likewise it can reach level $5q$ from above only by an active step taken backwards.

\subsection{The lens}

\begin{lemma}[the lens bound]\label{lem:lens}
Let $w,w'\in\C$ with $|w|,|w'|\le r$ and $w-w'=2e$, $|e|=1$. Let $U\in\C$, $C=|\re(Ue)|$,
$D=|\im(Ue)|$, and suppose $C\sqrt{r^2-1}\le D$. Then $|\re(U(w+w'))|\le 2D\sqrt{r^2-1}$.
\end{lemma}

\begin{proof}
Put $\xi=\tfrac12(w+w')\bar e=x+iy$. Then $w\bar e=\xi+1$ and $w'\bar e=\xi-1$, so $|\xi\pm1|\le r$:
the point $\xi$ lies in the lens of two disks of radius $r$ about $\pm1$, that is
$(|x|+1)^2+y^2\le r^2$. Now $\re(U(w+w'))=2\re(Ue\,\xi)$ is at most $2(C|x|+D|y|)$ in absolute
value, and on the lens $C|x|+D|y|\le D\sqrt{r^2-1}$: the function
$x\mapsto Cx+D\sqrt{r^2-(1+x)^2}$ is concave on $[0,r-1]$ with non-positive slope at $0$, so its
maximum is taken at the tips $x=0$, $y=\pm\sqrt{r^2-1}$.
\end{proof}

Apply this to an active step $v\to v+e$ whose ends both lie in the window, with $w=p+1-2v$,
$w'=p+1-2(v+e)$ and $U=\bar u$ (up to the rotation $m$, $w$ and $w'$ are the vectors from the two
centres to $p$). Since $\re(\bar u\cdot 2v)=\psi(v)$ and $\psi(e)=3-2\varphi$, and since both active
$e$ give $C=\varphi-\tfrac32$ and $D=\tfrac12\sqrt{7-4\varphi}$, for which the hypothesis of
\cref{lem:lens} holds when $r^2<3+\varphi$, we get
\begin{equation}\label{eq:window}
  |\psi(v)-Z_p|\le L(r) := \tfrac12\sqrt{7-4\varphi}\,\sqrt{r^2-1},
  \qquad Z_p=\re(\bar u(p+1))-\tfrac12(3-2\varphi).
\end{equation}
And $2L(r)<\delta$ if and only if $r^2<3+\varphi$: indeed $(2L)^2=(7-4\varphi)(r^2-1)$,
$\delta^2=(3-\varphi)^2$, and $(3-\varphi)^2=(7-4\varphi)(2+\varphi)$.

\subsection{Cut levels}

By \cref{lem:level}(i) the values of $\psi$ at the pivots on level $5q$ lie in the progression
$-5q+\delta\Z$, the \emph{teeth} of \cref{fig:levels}. Call $q$ a \emph{cut level} (for $p$) if this
progression misses the interval $[Z_p-L(r),Z_p+L(r)]$. By \eqref{eq:window}, no active step starts
on a cut level.

\begin{lemma}[cut levels]\label{lem:cut}
Let $\delta>0$, let $t\in\R$ with $t/\delta\notin\Q$, and let $0\le 2L<\delta$. There is $N$ such
that for every $Z\in\R$ some $q\in[1,N]$ and some $q\in[-N,-1]$ satisfy
$(tq+\delta\Z)\cap[Z-L,Z+L]=\varnothing$.
\end{lemma}

\begin{proof}
Divide by $\delta$: with $\theta=t/\delta$ and $\eta=L/\delta<\tfrac12$ the point
$q\theta-Z/\delta$ of $\R/\Z$ must enter an arc of length $1-2\eta>0$. By Dirichlet's approximation
theorem there are $a\ge1$ and $b\in\Z$ with $0<|a\theta-b|<1-2\eta$ ($\ne0$ as $\theta\notin\Q$).
Stepping $q$ by $a$ moves the point by the fixed small amount $a\theta-b$, which cannot jump over
the arc; starting at $q=\pm1$, it enters the arc within $\lceil 1/|a\theta-b|\rceil+1$ steps.
\end{proof}

We use it with $t=-5$: consecutive rows of teeth are shifted by $5\equiv\delta/\varphi\pmod\delta$,
and $5/\delta=2+\varphi$ is irrational --- a golden rotation.

\begin{figure}[t]
  \centering
  \includegraphics[width=0.66\textwidth]{figures/b2_levels.png}
  \caption{Levels of the orbit of \cref{fig:lattice} ($r=2$). Vertical: the level $\ell$; horizontal:
    $\psi$. On each row $5q$ the teeth mark the progression $-5q+\delta\Z$; the green strip is
    $[Z_p-L(r),Z_p+L(r)]$. Rows whose teeth all miss the strip are cut levels (red): no active step
    leaves them upwards, so the walk's pivots (blue) and tips (orange) stay in the band between the
    two nearest cut levels.}
  \label{fig:levels}
\end{figure}

\subsection{Confinement}

\begin{proposition}\label{prop:lower}
For every $r<\rs$ there is $N(r)$ such that every orbit of $\GG_5(r)$ has at most $N(r)$ points.
Hence $\GG_5(r)$ is finite.
\end{proposition}

\begin{proof}
Besides $\psi$ use the rotated functional $\psi_2(v)=\psi(\zeta^{-1}v)=2\langle\zeta u,v\rangle$, with
level $\ell_2(v)=\ell(\zeta^{-1}v)$; since $\zeta^{-1}$ preserves $\mu_5$, classes and window
distances, everything above holds for it, with its own centre $Z'_p$. As $2L(r)<\delta$,
\cref{lem:cut} gives $N$, depending on $r$ only, and cut levels $q_-<0<q_+$ for $(\psi,Z_p)$ and
$q'_-<0<q'_+$ for $(\psi_2,Z'_p)$, all in $[-N,N]$. Let
\[
  V=\{v\in\Win_p : \cls(v)\in\{0,1\},\ 5q_-<\ell(v)\le 5q_+,\ 5q'_-<\ell_2(v)\le 5q'_+\}.
\]
\emph{$V$ is closed.} A typed step from $v\in V$ into the window lands in class $0$ or $1$. It could
leave the first band only by an active step from a pivot on level $5q_+$ or onto a pivot on
level $5q_-$, with both ends in $\Win_p$; but $q_\pm$ are cut levels. The same holds for $\ell_2$.
The dumbbell $(0,1)$ has levels $(0,3)$ and $(0,-2)$, inside both bands, and one of its ends lies
in $\Win_p$ because $p$ lies in a disk; so \cref{lem:invariant} applies.

\emph{$V$ is small.} For $v\in V$ we have $|v|\le1+r$, hence $|\psi(v)|,|\psi_2(v)|\le 2+2r$, and
$|\ell(v)|,|\ell_2(v)|\le 5N$; by \cref{lem:level}(i), $\delta|B(v)|\le|\psi(v)|+|\ell(v)|$, and
likewise for $B_2$. The four integers $\ell,B,\ell_2,B_2$ determine $v$: in a basis of $\Lam$ they are
given by an integer matrix of determinant $-5$. So $V$ lies in a box of lattice points whose size
depends on $r$ only, and by \cref{lem:invariant} $\Or_r(p)$ has at most $5$ times as many points
as the box (the pivot, or the tip, lies in $V$, and a tip differs from its pivot by a unit). The
group is finite by \cref{prop:bounded}.
\end{proof}

\Cref{prop:upper,prop:lower} are parts (a) and (b) of \cref{thm:main}.

\subsection{Why the same number twice}

The two halves meet at the same radius for visible reasons (\cref{fig:threshold}). The words of
\cref{lem:words} act on the whole chord only if $E'E$ lies in the lens, that is $|E+1|\le r$; the
cut needs the window $2L(r)$ shorter than the period $\delta$. Both conditions switch at
$r^2=3+\varphi$. In both halves the decisive points are the tips of the lens: the maximum in
\cref{lem:lens} is attained there, and at $r=\rs$ the tips are the lattice points
$\pm(\zeta-\zeta^{-1})$, where the three words touch the rims.

\begin{figure}[t]
  \centering
  \includegraphics[width=0.9\textwidth]{figures/b3_threshold.png}
  \caption{At $r=\rs$: the chord $E'E$ fits the lens exactly (left), and the window $2L(r)$ is
    exactly one period wide, so every row of teeth meets it and no cut level is left (right, rows
    $q=-25,\dots,25$). For $r<\rs$ the strip is narrower and cut levels appear on both sides.}
  \label{fig:threshold}
\end{figure}

% =================================================================================================
\section{The formal proof}\label{sec:formal}

\Cref{thm:main} is proved in Lean~4~\cite{lean4} (version~4.34.0) with Mathlib~\cite{mathlib} (a
pinned commit) by the library \texttt{CriticalRadiusFive}, about $2\,000$ lines in eleven modules
(\cref{tab:modules}). The formal statements are those quoted in \cref{sec:result}; their proofs use
the axioms \texttt{propext}, \texttt{Classical.choice} and \texttt{Quot.sound} and nothing else, as
the file \texttt{Test.lean} checks with \verb|#print axioms|.

\begin{table}[h]
  \centering\small
  \begin{tabularx}{\textwidth}{@{}lX@{}}
    \toprule
    Module & Content \\
    \midrule
    \texttt{Defs} & the turns, $\GG_n(r)$, the critical radius (general $n$) \\
    \texttt{Basic} & how the turns act; orbits of points off both disks \\
    \texttt{Finiteness} & \cref{prop:bounded,prop:mono}; the critical radius as a threshold \\
    \texttt{Words} & words acting on segments; convexity (proof of \cref{lem:words}) \\
    \texttt{Cyclotomic} & exact arithmetic in $\Z[\zeta]$, classes modulo $1-\zeta$, the norm \\
    \texttt{CircleRotation} & irrational rotations: no return, and \cref{lem:cut} \\
    \texttt{UpperBound} & \cref{lem:words}, \cref{prop:upper} \\
    \texttt{Lens} & \cref{lem:lens} \\
    \texttt{PivotWalk} & the normal form \eqref{eq:normal} and \cref{lem:invariant} \\
    \texttt{LowerBound} & \cref{lem:level}, \eqref{eq:window}, \cref{prop:lower} \\
    \texttt{Main} & \cref{thm:main} \\
    \bottomrule
  \end{tabularx}
  \caption{The modules of the formal proof.}
  \label{tab:modules}
\end{table}

\paragraph{What exactly is checked.}
The file \texttt{Challenge.lean} imports only Mathlib. It contains the five definitions of
\texttt{Defs} (the turn of a disk, the two generators, the group and the critical radius) and the
two theorems of \cref{sec:result} with \texttt{sorry} in place of proofs; a reader who wants to know
what has been proved needs to read only this file. The file \texttt{Solution.lean} states the same
two theorems and proves them from the library. The comparator~\cite{comparator}, configured by
\texttt{comparator.json}, checks that each of the two theorems is proved in \texttt{Solution.lean}
with exactly the statement it has in \texttt{Challenge.lean} --- every constant the statement
mentions, definitions included, must be identical in the two modules --- that the proofs use only
the three axioms above, and that both Lean's kernel and the independent kernel \texttt{nanoda}
accept the solution; on Linux it runs every build and replay in a sandbox.
\placeholder{Confirm the comparator run and its output before publication.}

\paragraph{Scope.}
Every numbered statement of \cref{sec:result,sec:upper,sec:lower} and the two displayed equations
are formalised (\cref{app:lean}); the uniform orbit bound of \cref{thm:main}(b) is proved inside
\leanname{finite\_GG\_five\_of\_lt}. Not formalised, because nothing depends on them: the remark
in \cref{sec:upper}, the density of $\Lam$, and the discussion in the last subsection of
\cref{sec:lower}. The figures were computed in floating point and illustrate; they prove nothing.

\paragraph{Reproducing.}
\texttt{lake exe cache get \&\& lake build} builds the library, the challenge, the solution and the
axiom audit; \texttt{scripts/verify-comparator.sh} runs the comparator. The figures are the final
frames of the animations in \texttt{animations/}, rendered by \texttt{animations/render.sh}.

% =================================================================================================
\appendix
\section{From the paper to the Lean declarations}\label{app:lean}

All names live in the namespace \texttt{CriticalRadiusFive}; \texttt{Cyc} is $\Lam$ with
coordinates in the basis $1,\zeta,\zeta^2,\zeta^3$.

{\small
\begin{longtable}{@{}>{\raggedright\arraybackslash}p{0.29\textwidth}>{\raggedright\arraybackslash}p{0.67\textwidth}@{}}
  \toprule
  Paper & Lean \\
  \midrule
  \endfirsthead
  \toprule
  Paper & Lean \\
  \midrule
  \endhead
  \bottomrule
  \endfoot
  $a$, $b$, $\GG_5(r)$, $\rc(5)$ & \leanname{turn}, \leanname{genA}, \leanname{genB}, \leanname{GG}, \leanname{criticalRadius} \\
  \Cref{prop:mono} & \leanname{orbit\_subset\_orbit\_of\_le} \\
  \Cref{prop:bounded} & \leanname{finite\_of\_encard\_orbit\_le} \\
  consequences in \cref{sec:result} & \leanname{finite\_GG\_of\_le}, \leanname{infinite\_of\_infinite\_orbit}, \leanname{finite\_GG\_iff\_of\_forall\_lt}, \leanname{criticalRadius\_eq\_of\_forall\_lt} \\
  words on segments, convexity & \leanname{wordPerm}, \leanname{admissible\_of\_mem\_segment}, \leanname{wordPerm\_apply\_of\_mem\_segment} \\
  $E$; exact disk tests & \leanname{chordE}; \leanname{inDisk}, \leanname{admissibleZ}, \leanname{movesZ} \\
  \Cref{lem:words} & \leanname{piece₁}, \leanname{piece₂}, \leanname{piece₃} \\
  the golden rotation of the chord & \leanname{chordRot}, \leanname{exists\_mem\_GG\_apply\_smul\_chordE} \\
  no return of an irrational rotation & \leanname{iterate\_injective\_of\_irrational}, \leanname{infinite\_orbit\_of\_irrational\_rotation} \\
  \Cref{prop:upper} & \leanname{infinite\_orbit\_zero}, \leanname{infinite\_GG\_five} \\
  $\Lam$, classes, units & \leanname{Cyc}, \leanname{Cyc.cls}, \leanname{Cyc.units}, \leanname{Cyc.cls\_of\_mem\_units} \\
  window, closed sets & \leanname{Window}, \leanname{EdgeClosed} \\
  \eqref{eq:normal}, \cref{lem:invariant} & \leanname{Good}, \leanname{reach}, \leanname{orbit\_subset\_reach} \\
  $u$, $\psi$, $\ell$ & \leanname{cutU}, \leanname{psi}, \leanname{psi\_eq}, \leanname{level} \\
  \Cref{lem:level} (i), (ii), (iii) & \leanname{psi\_add\_level}; \leanname{level\_of\_mem\_units}, \leanname{psi\_of\_level\_eq\_three}; \leanname{level\_emod\_five} \\
  \Cref{lem:lens} & \leanname{abs\_re\_mul\_add\_le} (with \leanname{lens\_linear\_le}) \\
  \eqref{eq:window}; $2L(r)<\delta$ & \leanname{abs\_psi\_sub\_cutCenter\_le}; \leanname{two\_mul\_lensWidth\_lt} \\
  \Cref{lem:cut}; $5/\delta\notin\Q$ & \leanname{exists\_cut\_levels}; \leanname{irrational\_neg\_five\_div} \\
  cut levels of $\psi$ & \leanname{IsCutLevel}, \leanname{exists\_isCutLevel} \\
  no step crosses a cut level & \leanname{level\_ne\_of\_isCutLevel}, \leanname{level\_mem\_band} \\
  the set $V$; $V$ is closed & \leanname{cutSet}; \leanname{edgeClosed\_cutSet}, \leanname{good\_cutSet} \\
  $V$ is small & \leanname{abs\_psiGold\_le}, \leanname{mem\_box\_of\_abs\_le}, \leanname{box}, \leanname{mem\_box\_of\_mem\_cutSet} \\
  \Cref{prop:lower} & \leanname{finite\_GG\_five\_of\_lt} \\
  \Cref{thm:main} & \leanname{finite\_GG\_five\_iff}, \leanname{criticalRadius\_five} \\
\end{longtable}
}

% =================================================================================================
\begin{thebibliography}{9}
\bibitem{hkrsv} R.~A. Hearn, W.~Kretschmer, T.~Rokicki, B.~Streeter, E.~Vergo,
  \emph{Two-disk compound symmetry groups}, arXiv:2302.12950, 2023.
\bibitem{lean4} L.~de~Moura, S.~Ullrich, \emph{The Lean~4 theorem prover and programming language},
  in: Automated Deduction -- CADE 28, LNCS 12699, Springer, 2021, 625--635.
\bibitem{mathlib} The mathlib Community, \emph{The Lean mathematical library}, in: Proceedings of
  CPP 2020, ACM, 2020, 367--381.
\bibitem{comparator} \emph{comparator}: checking a Lean solution against a challenge,
  \url{https://github.com/leanprover/comparator}.
\end{thebibliography}

\end{document}
